\chapter{Probabilités conditionnelles et indépendance}
\thispagestyle{fancy}

\section{Probabilité conditionnelle}

\begin{definition}
    Soit $A$ et $B$ deux événements avec $P(A)\neq 0$.\\
    On appelle probabilité conditionnelle de $B$ sachant $A$, la probabilité que l'événement $B$ se réalise sachant que l'événement $A$ est réalisé. Elle est notée $P_A (B)$  et est définie par : \\$P_A(B)= \dfrac{Card(A \cap B)}{Card(A)}$. 
\end{definition}

\begin{bexemples}
%    \begin{enumEx}
    
%        \item	On tire une carte au hasard dans un jeu de 32 cartes.\\
%        Soit $A$ l'événement \guillemets{Le résultat est un pique}.\\
%        Soit $B$ l'événement \guillemets{Le résultat est un roi}.\\
%        Donc $A \cap B$ est l'événement \guillemets{Le résultat est le roi de pique}.\\
%        
%        Alors : $P(A)= \dfrac{8}{32} = \dfrac{1}{4}$ et $P(A \cap B)= \dfrac{1}{32}$.\\
%        Donc la probabilité que le résultat soit un roi sachant qu'on a tiré un pique est :\\
%        $P_A(B)= \dfrac{Card(A \cap B)}{Card(A)}=\dfrac{\dfrac{1}{32}}{\dfrac{1}{4}}=\dfrac{1}{8}$.\\
%        \ \\
%        On peut retrouver intuitivement ce résultat. En effet, sachant que le résultat est un pique, on a une chance sur 8 d'obtenir le roi.\\
%        \\
%        \item 
        Un sac contient 50 boules, dont 20 boules rouges et 30 boules noires, où il est marqué soit \guillemets{Gagné} ou soit \guillemets{Perdu}.\\ 
        Sur 15 boules rouges, il est marqué Gagné.\\
        Sur 9 boules noires, il est marqué Gagné.\\
        On tire au hasard une boule dans le sac.\\
        Soit $R$ l'événement \guillemets{On tire une boule rouge}.\\
        Soit $G$ l'événement \guillemets{On tire une boule marquée Gagné}.\\
        Donc $R \cap G$ est l'événement \guillemets{On tire une boule rouge marquée Gagné}.\\
        \ \\
        Alors : $P(R)= \dfrac{20}{50} = \dfrac{2}{5}=0,4$ et $P(R \cap G)= \dfrac{15}{50}=\dfrac{3}{10}=0,3$.\\
        Donc la probabilité qu'on tire une boule marquée Gagné sachant qu'elle est rouge est :\\
        $P_R(G)= \dfrac{Card(R \cap G)}{Card(R)}=\dfrac{O,3}{0,4}=\dfrac{3}{4}=0,75$.\\
        \ \\
        On peut retrouver intuitivement ce résultat. En effet, sachant que le résultat est une boule rouge, on a 15 chances sur 20 qu'il soit marqué Gagné.
%    \end{enumEx}
\end{bexemples}

\begin{brmq}
	La probabilité conditionnelle suit les règles et lois de probabilités vues pour les probabilité simples. On a en particulier :
\end{brmq}

\begin{bprop}
    Soit A et B deux événements avec $P(A)\neq 0$.
    \begin{itemProp}
    	\item $0\leq P_A(B)\leq 1$
    	\item $P_A(\overline{B})=1-P_A(B)$
    	\item $P(A \cap B)=P(A)\times P_A (B)$
    \end{itemProp}
\end{bprop}

\section{Arbre pondéré}

\subsection{Exemple}

\begin{bexemple}
    On reprend l'exemple étudié au paragraphe I.
    L'expérience aléatoire peut être schématisée par un arbre pondéré (ou arbre de probabilité) :
%    \begin{center}
%    	\includegraphics[scale=0.8]{01-09-01.png}
%    \end{center}
	\begin{center}
		%\includegraphics[scale=0.8]{TSTG.04.01.png}\\
		\begin{tikzpicture}[xscale=1,yscale=1]
			% Styles (MODIFIABLES)
			\tikzstyle{fleche}=[->,>=latex,thick]
			\tikzstyle{noeud}=[rectangle]
			\tikzstyle{feuille}=[rectangle]
			\tikzstyle{etiquette}=[midway,fill=white]
			% Dimensions (MODIFIABLES)
			\def\DistanceInterNiveaux{3.8}
			\def\DistanceInterFeuilles{2.5}
			% Dimensions calculées (NON MODIFIABLES)
			\def\NiveauA{(0)*\DistanceInterNiveaux}
			\def\NiveauB{(1.5)*\DistanceInterNiveaux}
			\def\NiveauC{(2.5)*\DistanceInterNiveaux}
			\def\InterFeuilles{(-1)*\DistanceInterFeuilles}
			% Noeuds (MODIFIABLES : Styles et Coefficients d'InterFeuilles)
			\node[noeud] (R) at ({\NiveauA},{(1.5)*\InterFeuilles}) {On tire une boule};
			\node[noeud] (Ra) at ({\NiveauB},{(0.5)*\InterFeuilles}) {Boule Rouge};
			\node[feuille] (Raa) at ({\NiveauC},{(0)*\InterFeuilles}) {Gagné };
			\node[feuille] (Rab) at ({\NiveauC},{(1)*\InterFeuilles}) {Perdu};
			\node[noeud] (Rb) at ({\NiveauB},{(2.5)*\InterFeuilles}) {Boule Noire};
			\node[feuille] (Rba) at ({\NiveauC},{(2)*\InterFeuilles}) {Gagné};
			\node[feuille] (Rbb) at ({\NiveauC},{(3)*\InterFeuilles}) {Perdu};
			% Arcs (MODIFIABLES : Styles)
			\draw[fleche] (R)--(Ra) node[etiquette] {\footnotesize{$P(R)=0,4$}};
			\draw[fleche] (Ra)--(Raa) node[etiquette] {\footnotesize{$P_R(G)=0,75$}};
			\draw[fleche] (Ra)--(Rab) node[etiquette] {\footnotesize{$P_R(\overline{G})$}};
			\draw[fleche] (R)--(Rb) node[etiquette] {\footnotesize{$P(\overline{R})=0,6$}};
			\draw[fleche] (Rb)--(Rba) node[etiquette] {\footnotesize{$P_{\overline{R}}(G)$}};
			\draw[fleche] (Rb)--(Rbb) node[etiquette] {\footnotesize{$P_{\overline{R}}(\overline{G})$}};
			\draw[ForestGreen,<-] (Raa)--($(Raa)+(1.5,0)$) node[right,font=\sffamily] {$R\cap G$ (boule rouge inscrit gagné)} ;
			\draw[white,<-] (Raa)--($(Raa)+(1.5,-0.5)$) node[right,font=\sffamily] {\textcolor{ForestGreen}{$P(R\cap G)=P(R)\times P(G)=0,3$}} ;
			\draw[ForestGreen,<-] (Rab)--($(Rab)+(1.5,0)$) node[right,font=\sffamily] {$R\cap \overline{G}$ (boule rouge inscrit perdu)} ;
			\draw[ForestGreen,<-] (Rba)--($(Rba)+(1.5,0)$) node[right,font=\sffamily] {$\overline{R}\cap G$ (boule noire inscrit gagné)} ;
			\draw[ForestGreen,<-] (Rbb)--($(Rbb)+(1.5,0)$) node[right,font=\sffamily] {$\overline{R}\cap \overline{G}$ (boule noire inscrit perdu)} ;
		\end{tikzpicture}
	\end{center}	
\end{bexemple}
\subsection{Règles}
\begin{bregle}
	La somme des probabilités des branches issues d'un même noeud est égale à 1.
\end{bregle}
\begin{bexemples}
    \begin{itemEx}
        \item A partir du noeud \guillemets{On tire une boule}, on a : $P(R)+P(\overline{R})=0,4+0,6=1$ 
        \item A partir du noeud \guillemets{Boule rouge}, on a : $P_R (\overline{G})=1-P_R(G)=1-0,75=0,25$.	
    \end{itemEx}
\end{bexemples}
Ces exemples font apparaître une formule donnée au paragraphe I.\\

\begin{bregle}
	La probabilité d'une \guillemets{feuille} (extrémité d'un chemin) est égale au produit des probabilités du chemin aboutissant à cette feuille.
\end{bregle}
\begin{bexemple}
	On considère la feuille $R \cap G$.
    On a :\\
     $P(R \cap G)=P(R)\times P_R (G)=0,4\times 0,75=0,3$ 
\end{bexemple}

\begin{bregle}
	\grasInProp{Formule des probabilités totales}\\
	La probabilité d'un événement associé à plusieurs "feuilles" est égale à la somme des probabilités de chacune de ces "feuilles".
\end{bregle}
\begin{bexemple}
    L'événement \guillemets{On tire une boule marquée Gagné} est associé aux feuilles $R \cap G$ et $\overline R \cap G$. On a :\\
    
    $P(R \cap G)=0,3$ et 
    $P(\overline R \cap G)= \dfrac{9}{50} = 0,18$\\ (Probabilité de tirer une boule noire marquée Gagné)\\
    Donc $P(G)=P(R \cap G)+P(\overline R \cap G)=0,3+0,18=0,48$.
\end{bexemple}

%\begin{bmethodet}{Calculer la probabilité d'un événement associé à plusieurs feuilles}
%    Lors d'une épidémie chez des bovins, on s'est aperçu que si la maladie est diagnostiquée suffisamment tôt chez un animal, on peut le guérir ; sinon la maladie est mortelle.\\ 
%    Un test est mis au point et essayé sur un échantillon d'animaux dont 2 \% est porteur de la maladie. On obtient les résultats suivants :\\
%    – si un animal est porteur de la maladie, le test est positif dans 85 \% des cas ;\\
%    – si un animal est sain, le test est négatif dans 95 \% des cas.\\
%    \ \\
%    On choisit de prendre ces fréquences observées comme probabilités pour toute la population et d'utiliser le test pour un dépistage préventif de la maladie.\\
%    \ \\
%    On note respectivement $M$ et $T$ les événements \guillemets{Être porteur de la maladie} et
%    \guillemets{Avoir un test positif}.
%    
%    1) Un animal est choisi au hasard. Quelle est la probabilité que son test soit positif ?\\
%    \emph{D'après BAC S, Antilles-Guyanne 2010}\\	
%    \ \\
%    2) Si le test du bovin est positif, quelle est la probabilité qu'il soit malade ?\\
%    \ \\
%    \ \\
%%    \begin{center}
%%    	\includegraphics{01-09-02}
%%    \end{center}
%		    \begin{center}
%		%\includegraphics[scale=0.8]{TSTG.04.01.png}\\
%		\begin{tikzpicture}[xscale=1,yscale=1]
%			% Styles (MODIFIABLES)
%			\tikzstyle{fleche}=[->,>=latex,thick]
%			\tikzstyle{noeud}=[rectangle]
%			\tikzstyle{feuille}=[rectangle]
%			\tikzstyle{etiquette}=[midway,fill=white]
%			% Dimensions (MODIFIABLES)
%			\def\DistanceInterNiveaux{4.5}
%			\def\DistanceInterFeuilles{1.5}
%			% Dimensions calculées (NON MODIFIABLES)
%			\def\NiveauA{(0)*\DistanceInterNiveaux}
%			\def\NiveauB{(1.5)*\DistanceInterNiveaux}
%			\def\NiveauC{(2.5)*\DistanceInterNiveaux}
%			\def\InterFeuilles{(-1)*\DistanceInterFeuilles}
%			% Noeuds (MODIFIABLES : Styles et Coefficients d'InterFeuilles)
%			\node[noeud] (R) at ({\NiveauA},{(1.5)*\InterFeuilles}){};
%			\node[noeud] (Ra) at ({\NiveauB},{(0.5)*\InterFeuilles}) {$M$};
%			\node[feuille] (Raa) at ({\NiveauC},{(0)*\InterFeuilles}) {$T$};
%			\node[feuille] (Rab) at ({\NiveauC},{(1)*\InterFeuilles}) {$\overline{T}$};
%			\node[noeud] (Rb) at ({\NiveauB},{(2.5)*\InterFeuilles}) {$\overline{M}$};
%			\node[feuille] (Rba) at ({\NiveauC},{(2)*\InterFeuilles}) {$T$};
%			\node[feuille] (Rbb) at ({\NiveauC},{(3)*\InterFeuilles}) {$\overline{T}$};
%			% Arcs (MODIFIABLES : Styles)
%			\draw[fleche] (R)--(Ra) node[etiquette] {\footnotesize{$P(M)=0,02$}};
%			\draw[fleche] (Ra)--(Raa) node[etiquette] {\footnotesize{$P_M(T)=0,85$}};
%			\draw[fleche] (Ra)--(Rab) node[etiquette] {\footnotesize{$P_M(\overline{T})=0,15$}};
%			\draw[fleche] (R)--(Rb) node[etiquette] {\footnotesize{$P(\overline{M})=0,98$}};
%			\draw[fleche] (Rb)--(Rba) node[etiquette] {\footnotesize{$P_{\overline{M}}(T)=0,05$}};
%			\draw[fleche] (Rb)--(Rbb) node[etiquette] {\footnotesize{$P_{\overline{M}}(\overline{T})=0,95$}};
%			\draw[ForestGreen,<-] (Raa)--($(Raa)+(1.5,0)$) node[right,font=\sffamily] {$M\cap T$} ;
%			%\draw[white,<-] (Raa)--($(Raa)+(1.5,-0.5)$) node[right,font=\sffamily] {\textcolor{ForestGreen}{$P(R\cap G)=P(R)\times P(G)=0,3$}} ;
%			\draw[ForestGreen,<-] (Rab)--($(Rab)+(1.5,0)$) node[right,font=\sffamily] {$M\cap \overline{T}$} ;
%			\draw[ForestGreen,<-] (Rba)--($(Rba)+(1.5,0)$) node[right,font=\sffamily] {$\overline{M}\cap T$} ;
%			\draw[ForestGreen,<-] (Rbb)--($(Rbb)+(1.5,0)$) node[right,font=\sffamily] {$\overline{M}\cap \overline{T}$} ;
%		\end{tikzpicture}
%	\end{center}	
%    La probabilité que le test soit positif est associée aux deux feuilles $M\cap T$ et $\overline M\cap T$.\\
%    $P(T)=P(M\cap T)+P(\overline M\cap T )$ (Formule des probabilités totales)\\
%             $P(T)= 0,02 \times 0,85 + 0,98 \times 0,05 = 0,066$.
%    La probabilité que le test soit positif est égale à 6,6\%. \\
%    \ \\
%    2) $P_T (M)= \dfrac{P(T\cap M)}{P(T)}  = \dfrac{0,02\times 0,85}{0,066} \approx 0,26$.\\
%    La probabilité que le bovin soit malade sachant que le test est positif est d'environ 26\%.
%\end{bmethodet}

\section{Probabilités et indépendance}

\subsection{Indépendance de deux événements}

\begin{bdefi}
	On dit que deux évènements $A$ et $B$ de probabilité non nulle sont indépendants lorsque $P(A\cap B)=P(A)\times P(B)$.
\end{bdefi}

\begin{brmq}
	On a également :
$A$ et $B$ sont indépendants, si et seulement si, $P_A (B)=P(B)$ ou $P_B (A)=P(A)$.
\end{brmq}

\begin{bexemple}
	On tire une carte au hasard dans un jeu de 32 cartes.\\
    Soit $R$ l'événement \guillemets{On tire un roi}.\\
    Soit $T$ l'événement \guillemets{On tire un trèfle}.\\
    Alors $R\cap T$ est l'événement \guillemets{On tire le roi de trèfle}.\\
    \ \\
    On a :\\
    $P(R)  = \dfrac{4}{32} = \dfrac{1}{8}$ , $P(T)= \dfrac{8}{32} = \dfrac{1}{4}$ et $P(R\cap T)= \dfrac{1}{32}$\\ 
    Donc $P(R)\times P(T)  = \dfrac{1}{8} \times \dfrac{1}{4} = \dfrac{1}{32} = P(R\cap T)$\\
    Les événements $R$ et $T$ sont donc indépendants.
    Ainsi, par exemple, $P_T (R)=P(R)$. Ce qui se traduit par la probabilité de tirer un roi parmi les trèfles et égale à la probabilité de tirer un roi parmi toutes les cartes.\\
    \\
    \underline{Contre-exemple :}\\
    On reprend l'expérience précédente en ajoutant deux jokers au jeu de cartes.\\
    Ainsi :
    $P(R)  = \dfrac{4}{34} = \dfrac{2}{17}$, $P(T)= \dfrac{8}{34} = \dfrac{4}{17}$ et $P(R\cap T)= \dfrac{1}{34}$\\
    Donc $P(R)\times P(T)  = \dfrac{2}{17} \times \dfrac{4}{17} = \dfrac{8}{289} \neq P(R\cap T)$\\
    Les événements R et T ne sont donc pas indépendants.
\end{bexemple}

%\begin{bmethodet}{Utiliser l'indépendance de deux événements}
%    Dans une population, un individu est atteint par la maladie $m$ avec une probabilité égale à $0,005$ et par la maladie $n$ avec une probabilité égale à $0,01$.\\
%    On choisit au hasard un individu de cette population.\\
%    Soit $M$ l'événement \guillemets{L'individu a la maladie $m$}.\\
%    Soit $N$ l'événement \guillemets{L'individu a la maladie $n$}.\\
%    On suppose que les événements $M$ et $N$ sont indépendants.\\
%    Calculer la probabilité de l'événement $E$ \guillemets{L'individu a au moins une des deux maladies}.\\
%    \\
%    $P(E)=P(M\cup N)=P(M)+P(N)-P(M\cap N)$\\
%             $P(E)=P(M)+P(N)-P(M)\times P(N)$, car les événements $M$ et $N$ sont indépendants.\\
%             $P(E)= 0,005 + 0,01 – 0,005 \times 0,01 = 0,01495$\\
%    
%    La probabilité qu'un individu choisi au hasard ait au moins une des deux maladies est égale à $0,01495$.\\
%\end{bmethodet}

\begin{bprop}
	Si $A$ et $B$ sont indépendants alors $\overline{A}$ et $B$ sont indépendants.
\end{bprop}

%    \begin{bdemo}
%    $P(\overline{A}\cap B)=P(B\cap \overline{A})$\\
%    $P(\overline{A}\cap B)=P(B)\times P_B (\overline{A})$\\
%    $P(\overline{A}\cap B)=P(B)\times (1-P_B(A))$\\
%    $P(\overline{A}\cap B)=P(B)\times (1-P(A))$ car A et B sont indépendants\\
%    $P(\overline{A}\cap B)=P(B)\times P(\overline{A})$\\
%    \\
%    Donc $\overline{A}$ et $B$ sont indépendants.
%\end{bdemo}

\begin{bexemple}
    Lors d'un week-end prolongé, Bison futé annonce qu'il y a $42\%$ de risque de tomber dans un bouchon sur l'autoroute A6 et $63\%$ sur l'autoroute A7.\\
    Soit $A$ l'événement \guillemets{On tombe dans un bouchon sur l'autoroute A6.}\\
    Soit $B$ l'événement \guillemets{On tombe dans un bouchon sur l'autoroute A7.}\\
    On suppose que les événements $A$ et $B$ sont indépendants.\\
    \\
    Alors les événements $\overline{A}$ et $B$ sont également indépendants et on a :\\
    $P(\overline{A} \cap B)=P(\overline{A})\times P(B)=0,58\times 0,63=0,3654$\\
    \\
    On peut interpréter ce résultat : \\
    La probabilité de tomber dans un bouchon sur l'autoroute A7 mais pas sur l'autoroute A6 est égale à $36,54\%$.
\end{bexemple}

\subsection{Succession de deux épreuves indépendantes}

\begin{bexemples}
    \begin{itemEx}
        \item On lance un dé plusieurs fois de suite et on note à chaque fois le résultat. On répète ainsi la même expérience (lancer un dé) et les expériences sont indépendantes l'une de l'autre (un lancer n'influence pas le résultat d'un autre lancer).
        \item Une urne contient 2 boules blanches et 3 boules noires. On tire au hasard une boule et on la remet dans l'urne.\\
        On répète cette expérience 10 fois de suite. Ces expériences sont identiques et indépendantes.
    \end{itemEx}
\end{bexemples}

\begin{bdefi}
	Plusieurs expériences sont \grasInDef{identiques et indépendantes} si :
    \begin{itemDef}
        \item elles ont les mêmes issues,
        \item chaque issue possède la même probabilité.
    \end{itemDef} 
\end{bdefi}

\begin{bprop}
	On considère une expérience aléatoire à deux issues $A$ et $B$ avec les probabilités $P(A)$ et $P(B)$.\\
    Si on répète l'expérience deux fois de suite :\\
    \begin{itemDef}
        \item la probabilité d'obtenir l'issue $A$ suivie de l'issue $B$ est égale à $P(A) \times P(B)$,
        \item la probabilité d'obtenir l'issue $B$ suivie de l'issue $A$ est égale à $P(B) \times P(A)$,
        \item la probabilité d'obtenir deux fois l'issue $A$ est égale à $P(A)^2$,
        \item la probabilité d'obtenir deux fois l'issue $B$ est égale à $P(B)^2$.
    \end{itemDef}
\end{bprop}

%\begin{bmethodet}{Représenter la répétition d'expériences identiques et indépendantes dans un arbre}
%On considère l'expérience suivante :\\
%Une urne contient 3 boules blanches et 2 boules rouges. On tire au hasard une boule et on la remet dans l'urne. On répète l'expérience deux fois de suite.\\
%1) Représenter l'ensemble des issues de ces expériences dans un arbre.\\
%2) Déterminer la probabilité :\\
%\ \ \ \ \ \ a) d'obtenir deux boules blanches\\
%\ \ \ \ \ \ b) une boule blanche et une boule rouge\\
%\ \ \ \ \ \ c) au moins une boule blanche.\\
%\\
%1) On note $A$ l'issue "On tire une boule blanche" et $B$ l'issue "On tire une boule rouge".
%$P(A) = 3/5 = 0,6$ et $P(B) = 2/5 = 0,4$.\\
%On résume les issues de l'expérience dans un arbre de probabilité :\\
%%\begin{center}
%%	\includegraphics[scale=0.7]{01-09-03}
%%\end{center}
%		    \begin{center}
%	%\includegraphics[scale=0.8]{TSTG.04.01.png}\\
%	\begin{tikzpicture}[xscale=1,yscale=1]
%		% Styles (MODIFIABLES)
%		\tikzstyle{fleche}=[->,>=latex,thick]
%		\tikzstyle{noeud}=[rectangle]
%		\tikzstyle{feuille}=[rectangle]
%		\tikzstyle{etiquette}=[midway,fill=white]
%		% Dimensions (MODIFIABLES)
%		\def\DistanceInterNiveaux{3}
%		\def\DistanceInterFeuilles{1.5}
%		% Dimensions calculées (NON MODIFIABLES)
%		\def\NiveauA{(0)*\DistanceInterNiveaux}
%		\def\NiveauB{(1.5)*\DistanceInterNiveaux}
%		\def\NiveauC{(2.5)*\DistanceInterNiveaux}
%		\def\InterFeuilles{(-1)*\DistanceInterFeuilles}
%		% Noeuds (MODIFIABLES : Styles et Coefficients d'InterFeuilles)
%		\node[noeud] (R) at ({\NiveauA},{(1.5)*\InterFeuilles}){};
%		\node[noeud] (Ra) at ({\NiveauB},{(0.5)*\InterFeuilles}) {$A$};
%		\node[feuille] (Raa) at ({\NiveauC},{(0)*\InterFeuilles}) {$A$};
%		\node[feuille] (Rab) at ({\NiveauC},{(1)*\InterFeuilles}) {$B$};
%		\node[noeud] (Rb) at ({\NiveauB},{(2.5)*\InterFeuilles}) {$B$};
%		\node[feuille] (Rba) at ({\NiveauC},{(2)*\InterFeuilles}) {$A$};
%		\node[feuille] (Rbb) at ({\NiveauC},{(3)*\InterFeuilles}) {$B$};
%		% Arcs (MODIFIABLES : Styles)
%		\draw[fleche] (R)--(Ra) node[etiquette] {\footnotesize{$0,6$}};
%		\draw[fleche] (Ra)--(Raa) node[etiquette] {\footnotesize{$0,6$}};
%		\draw[fleche] (Ra)--(Rab) node[etiquette] {\footnotesize{$0,4$}};
%		\draw[fleche] (R)--(Rb) node[etiquette] {\footnotesize{$0,4$}};
%		\draw[fleche] (Rb)--(Rba) node[etiquette] {\footnotesize{$0,6$}};
%		\draw[fleche] (Rb)--(Rbb) node[etiquette] {\footnotesize{$0,4$}};
%		\draw[ForestGreen,<-] (Raa)--($(Raa)+(1.5,0)$) node[right,font=\sffamily] {$0,6\times 0,6=0,36$} ;
%		%\draw[white,<-] (Raa)--($(Raa)+(1.5,-0.5)$) node[right,font=\sffamily] {\textcolor{ForestGreen}{$P(R\cap G)=P(R)\times P(G)=0,3$}} ;
%		\draw[ForestGreen,<-] (Rab)--($(Rab)+(1.5,0)$) node[right,font=\sffamily] {$0,6\times 0,4=0,24$} ;
%		\draw[ForestGreen,<-] (Rba)--($(Rba)+(1.5,0)$) node[right,font=\sffamily] {$0,4\times 0,6=0,24$} ;
%		\draw[ForestGreen,<-] (Rbb)--($(Rbb)+(1.5,0)$) node[right,font=\sffamily] {$0,4\times 0,4=0,16$} ;
%	\end{tikzpicture}
%\end{center}	
%2)	a) Obtenir deux boules blanches correspond à l'issue $(A ; A)$ :\\
%	$P_1 =  0,36$ (d'après l'arbre).\\
%\\	
%	b) Obtenir une boule blanche et une boule rouge correspond aux issues 
%$(A ; B)$ et $(B ; A)$ :\\
%	$P_2 = 0,24 + 0,24 = 0,48$.\\
%\\
%	b) Obtenir au moins une boule blanche correspond aux issues 
%$(A ; B)$, $(A ; A)$ et $(B ; A)$ :\\
%	$P_3 = 0,24 + 0,36 + 0,24 = 0,84$.
%\end{bmethodet}

\begin{brmq}
    \begin{itemRem}
        \item Pour une expérience dont le nombre d'issues est supérieur à 2, le principe reste le même.
        \item Pour une expérience dont le nombre de répétition est supérieur à 2, le principe reste le même.
    \end{itemRem}
\end{brmq}

\begin{bexemple}
    On lance un dé à six faces 4 fois de suite.\\
    On considère les issues suivantes :\\
    A : On obtient un nombre pair.\\
    B : On obtient un 1.\\
    C : On obtient un 3 ou un 5.\\
    La probabilité d'obtenir la suite d'issues $(A ; B ; A ; C)$ est égale à :\\
    $\dfrac{1}{2}\times \dfrac{1}{6} \times \dfrac{1}{2} \times \dfrac{1}{3}=\dfrac{1}{72}$
\end{bexemple}
